\section{Формирование требований к системе}

TODO вводная часть для связи с предыдущим разделом

\subsection*{Функциональные требования}

\begin{enumerate}

\item \textbf{Управление источниками и подписками:} Система должна предоставлять возможность добавления ссылок на веб-сайты и публичные Telegram-каналы для их фонового мониторинга и парсинга.

\item \textbf{Индивидуальная параметрическая фильтрация:} Пользователь должен иметь возможность настроить фильтрацию по городу, временному диапазону, формату (онлайн/офлайн), стоимости (платные/бесплатные) и набору категорий (примеры TODO).

\item \textbf{AI-суммаризация и стандартизация анонсов:} Модуль сбора данных должен пропускать неструктурированный текст через LLM для формирования стандартизированной карточки события, включающей разедыл:
Название, Место, Дата, Цена, Ссылка, Краткое описание. 
Тексты, не являющиеся анонсами IT-мероприятий, должны автоматически игнорироваться.

\item \textbf{Омниканальная маршрутизация:} Доставка дайджестов пользователям должна осуществляться в Telegram-боте и на электронную поочту пользователя на основе его настроек профиля.

\item \textbf{Экспорт данных анонсов:} Интеграция функции генерации валидных \texttt{.ics} файлов для добавления мероприятий в цифровой календарь пользователя.

\item \textbf{Веб-интерфейс:} Наличие интерактивного веб-клиента для авторизации через OAuth, управления профилем, связывания аккаунтов (Telegram и Web) и удобной настройки фильтров поиска анонсов.

\end{enumerate}

\subsection*{Нефункциональные требования}

\begin{enumerate}

\item \textbf{Архитектура:} Приложение должно быть разделено на независимые сервисы: Scrapper (ядро бизнес-логики, включающее модуль сбора данных), Bot (интерфейс Telegram) и Web Frontend (клиентский интерфейс). 
Контракты взаимодействия должны описываться в формате OpenAPI.

\item \textbf{Отказоустойчивость:} Реализация механизмов Circuit Breaker, таймаутов, Retry в интеграциях модудей LLM, Telegram API, SMTP. 
Использование паттерна Fallback для переключения между асинхронным брокером сообщений Kafka и HTTP.

\item \textbf{Производительность:} Использование асинхронного взаимодействия сервисов. Применение пакетной обработки (батчинга) с пагинацией для предотвращения утечек памяти при работе с БД.

\item \textbf{Наблюдаемость:} Экспорт бизнес- и системных метрик (RED-метрики, потребление памяти, статистика парсинга) в формате Prometheus для визуализации в Grafana.

\item \textbf{Целостность данных:} Валидация входных данных как на стороне клиентского интерфейса, так и на сервере; проверка доступности источников для парсинга при добавлении ссылок на них в систему.

\end{enumerate}
